\documentclass{situationist-preprint}

\title{Narrative Identity and Its Co-Authors}
\msid{PID-09-a}
\msfamily{PID-09}
\msdiscipline{Personal Identity and the Self}
\mstype{Principal manuscript}
\mskeywords{narrative identity; self-constitution; characterization; numerical identity; co-authorship; recognition; Schechtman; MacIntyre}

\begin{document}
\maketitle

\begin{abstract}
Narrative theorists hold that a person's identity is constituted by the story of her life; their critics reply that lives are not stories and that many people do not experience them as such. We argue that narrative identity is real, that it is co-authored by others, and that it is distinct from numerical identity. Narrative identity answers the characterization question, which actions, experiences and traits are a person's own and what they mean, not the reidentification question of what makes a later person the same as an earlier one. It is real because it figures in explanations of conduct, emotion and social practice that cannot be given without it. It is co-authored because a self-narrative governs a person's conduct and others' treatment of her only with their uptake, and because the meaning of past acts is fixed partly by later events, including the acts of others. It is distinct from numerical identity because it differs from it in properties: it comes in degrees, it is revisable and it is not settled by the subject alone, while numerical identity is none of these. We answer objections from episodic lives, from the literary disanalogy and from autonomy.
\end{abstract}

\section{The problem}

\citet[ch.~15]{MacIntyre1981} argued that the unity of a human life is the unity of a narrative quest, and that we are at most co-authors of the narratives we live. \citet{Ricoeur1992} distinguished sameness, the identity of a thing that remains the same, from selfhood, the identity of one who keeps her word and answers for her past, and held that narrative mediates between them. \citet{Schechtman1996} separated two questions that the literature on personal identity had run together: the \emph{reidentification question}, what makes a person at one time the same person as a person at another, and the \emph{characterization question}, which actions, experiences and characteristics are attributable to a person; and she argued that the second is answered by a narrative self-conception. \citet{GStrawson2004} replied that narrativity is neither a universal feature of human experience nor a condition of a good life, since some people, whom he calls episodic, do not experience themselves as having a past or future in any story-like way. \citet{Lamarque2004} argued that lives lack the closure, selection and interpretive structure that make literary narratives narratives.

The question whether narrative identity is real therefore has two parts: whether there is a phenomenon that the narrative theorist describes, and whether it is identity in the sense of PID-01. We argue that the first answer is yes, and the second no; and we add a thesis the critics have not addressed, that narrative identity is not the work of the subject alone.

\section{Framework}

\begin{definition}[Narrative identity]\label{d:ni}
A person's \emph{narrative identity} at a time is the structure of attributions, of actions, experiences, traits and commitments to her, together with their significance, organized by an interpretation of her life as a temporally extended whole, that is authoritative for her conduct and for others' treatment of her.
\end{definition}

\begin{definition}[Co-authorship]\label{d:co}
A structure of attributions is \emph{co-authored} if its content is fixed in part by the contributions of persons other than its subject, so that the subject cannot alone determine it.
\end{definition}

\begin{definition}[Numerical identity]\label{d:num}
\emph{Numerical identity} of persons is the relation that holds between a person at one time and a person at another when they are one and the same; by PID-01, its condition for persons is non-branching psychological continuity, with bodily continuity as a tie-breaker.
\end{definition}

\begin{principle}[Explanatory criterion]\label{p:crit}
A structure exists if it has instances and figures in true explanations that cannot be reformulated, without loss of explanatory power, in terms of other things.
\end{principle}

\begin{principle}[Uptake]\label{p:uptake}
A self-interpretation is authoritative for a person's conduct and for others' treatment of her only to the extent that the attributions it contains receive uptake from others: that they are acknowledged, relied on and responded to in shared practices of credit, blame, trust and expectation.
\end{principle}

Principle~\ref{p:crit} is the criterion used in MET-20-a and PID-04-a. Principle~\ref{p:uptake} is a claim about authority, not about belief: one may privately interpret oneself as one likes, but an interpretation has the role Definition~\ref{d:ni} requires only if it is the one on which one's commitments are relied upon, one's past acts credited or held against one, and one's conduct anticipated. A person who takes herself to be a reliable friend, while all her friends treat her as unreliable, has at best a disputed narrative identity as a reliable friend.

\section{The argument}

\begin{lemma}\label{l:real}
Narrative identity is real.
\end{lemma}

\begin{proof}
Every person who lives among others has a structure of the kind in Definition~\ref{d:ni}: there is a fact, at any time, about which past acts she is answerable for, which commitments she is taken to have and what her history is taken to mean, and these facts govern her conduct and others' responses to her. The structure figures in explanations that cannot be given without it. Regret and pride are explained by the significance an act has within an interpretation of the agent's life, not by the act alone; a person keeps a commitment made years before because it belongs to who she takes herself to be; a community's trust in her is explained by the history it attributes to her. Replacing the structure by a list of the person's past states loses the explanation, since the same states, differently organized and differently understood, explain different responses. By Principle~\ref{p:crit}, narrative identity exists. The conclusion agrees with PID-04, on which the self is a construction that includes the organism's self-interpretation, and constructions are real.
\end{proof}

\begin{lemma}\label{l:co}
Narrative identity is co-authored.
\end{lemma}

\begin{proof}
There are three contributions by others. First, by Principle~\ref{p:uptake}, the authority that makes a self-interpretation part of narrative identity depends on its uptake by others, and others can withhold, qualify or redirect it. Second, the materials of the narrative are partly supplied by others: a person's early life is known to her chiefly through the stories of those who raised her, and many of her actions are what they are, a promise kept or a betrayal, only in virtue of the expectations and responses of the persons involved. Third, the significance of a past act depends in part on what follows it. An insult that is later forgiven, or a failure that later becomes the beginning of a vocation, has a significance it did not have when it occurred; REL-22-a argues that significance properties of this kind are settled only as later events are settled, and that the free acts of others, such as forgiveness, are among the events that settle them. Since the content of narrative identity includes the significance of the attributed acts (Definition~\ref{d:ni}), others' later acts contribute to fixing it. By Definition~\ref{d:co}, narrative identity is co-authored.
\end{proof}

\begin{lemma}\label{l:dist}
Narrative identity is distinct from numerical identity.
\end{lemma}

\begin{proof}
If two relations are the same, they have the same properties. Narrative identity comes in degrees: a life can be more or less integrated by its interpretation, and an attribution more or less secure. It is revisable: by Lemma~\ref{l:co}, later events can change the significance of earlier acts, and hence the content of the structure, without changing which person did them. And it is not settled by the subject alone. Numerical identity, as Definition~\ref{d:num} specifies it, is none of these: whether a later person is the same as an earlier one is fixed by the continuity between them, it is not changed by later reinterpretation, and it does not depend on anyone's uptake. The extensions also diverge. A person who has lost all autobiographical memory survives (PID-08), while the narrative that organized her life is in large part lost to her and carried on by others; and a person may be credited, by herself and by others, with an act that was done by someone else, or fail to have integrated an act she did, without any change in who she is.
\end{proof}

\begin{theorem}\label{t:main}
Narrative identity is real; it is co-authored by others; and it is distinct from numerical identity.
\end{theorem}

\begin{proof}
By Lemmas~\ref{l:real}, \ref{l:co} and~\ref{l:dist}.
\end{proof}

Theorem~\ref{t:main} vindicates Schechtman's separation of the two questions while rejecting the view that the narrative answers both. It also gives MacIntyre's remark about co-authorship a precise sense: the subject is the principal author of her narrative identity, in that her own self-interpretation is the default content, but she is not its sole author, and in childhood, illness and disgrace others may hold the larger share.

\section{Objections and replies}

\begin{objection}[Episodic lives]
Some people do not experience their lives as stories and do not tell them to themselves \citep{GStrawson2004}. On the present view they lack narrative identity, which is implausible, or the view overstates the phenomenon.
\end{objection}

\begin{reply}
Definition~\ref{d:ni} does not require that the subject experience her life as a story, or tell one. It requires an interpretive structure that is authoritative for her conduct and for others' treatment of her. An episodic person keeps promises, is credited and blamed for past acts and is trusted or distrusted on the basis of a history; that structure exists whether or not she narrates it, and Lemma~\ref{l:co} explains why it can exist when its subject contributes little to it. Strawson's point refutes the claim that a good life requires self-narration, which we do not make.
\end{reply}

\begin{objection}[Lives are not stories]
Narratives have authors, plots and endings selected for meaning. Lives have none of these, so a life cannot have a narrative identity except by courtesy \citep{Lamarque2004}.
\end{objection}

\begin{reply}
The objection holds the phenomenon to literary standards it does not claim to meet. What Definition~\ref{d:ni} calls narrative is an interpretation of a life as a temporally extended whole in which earlier and later events bear on one another's significance. That structure need not have closure; Lemma~\ref{l:co} implies that it cannot have it while the life and its consequences continue. If \enquote{narrative} is reserved for literary form, the thesis can be restated with another word without loss.
\end{reply}

\begin{objection}[Autonomy]
If others co-author one's identity, one is at their mercy. Oppressive communities would then have the power to make people who they say they are.
\end{objection}

\begin{reply}
They have some such power, and the view explains it rather than inventing it. \citet{LindemannNelson2001} describes how dominant stories damage the identities of members of oppressed groups by constraining the uptake their self-interpretations receive, and how counterstories repair them by securing uptake in another community. Co-authorship is not dictation: the subject's interpretation is the default, uptake can be sought from more than one audience, and a disputed narrative identity is not the same as an imposed one. PID-12 argues more generally that the self is partly, not wholly, constituted by recognition.
\end{reply}

\section{What the result does not establish}

The result does not make numerical identity depend on narrative, and does not revise PID-01. It does not claim that narrative identity is what matters in survival; PID-06 argues that psychological connectedness carries most of what matters, and narrative integration is at most one form connectedness takes. Nor does it say how authorship is shared in particular cases, or that every contribution of others is legitimate.

\section{Conclusion}

There is a fact about who, in the characterizing sense, a person is: which acts are hers, what they mean and what she is answerable for. That fact is real, because it explains what nothing else explains; it is fixed in part by others, through uptake, through the materials they supply and through what they later do; and it differs from the fact that she is one person over time, which is fixed by continuity and is untouched by reinterpretation.

\begin{thebibliography}{99}\small
\bibitem[Lamarque(2004)]{Lamarque2004} Lamarque, P. (2004). On not expecting too much from narrative. \emph{Mind \& Language}, 19(4), 393--408.
\bibitem[Lindemann~Nelson(2001)]{LindemannNelson2001} Lindemann Nelson, H. (2001). \emph{Damaged Identities, Narrative Repair}. Cornell University Press.
\bibitem[MacIntyre(1981)]{MacIntyre1981} MacIntyre, A. (1981). \emph{After Virtue: A Study in Moral Theory}. University of Notre Dame Press.
\bibitem[Ricoeur(1992)]{Ricoeur1992} Ricoeur, P. (1992). \emph{Oneself as Another}, trans. K. Blamey. University of Chicago Press.
\bibitem[Schechtman(1996)]{Schechtman1996} Schechtman, M. (1996). \emph{The Constitution of Selves}. Cornell University Press.
\bibitem[Strawson(2004)]{GStrawson2004} Strawson, G. (2004). Against narrativity. \emph{Ratio}, 17(4), 428--452.
\end{thebibliography}

\end{document}
